
\[
  \underline{\underline{D}}_{\pi'}(\underline{E})\simeq \underline{\underline{D}}_{\pi''}(\underline{F}).
\]

La
démonstration, immédiate sur la définition fonctorielle des schémas de drapeaux, est laissée au lecteur.

\textbf{Lemme 4.3.} \uline{Soient} \(A\) \uline{un anneau commutatif avec unité}, \(n\) \uline{un entier} \(>0\), \((c_i)_{1\le i\le n}\) \uline{une famille d'éléments de} \(A\) ; \uline{soit} \(C\) \uline{la} \(A\)-\uline{algèbre engendrée par des générateurs} \(\xi_i\), \(i=1,\ldots,n\), \uline{soumis aux relations}~:
\[
  \sigma_i(\xi_1,\ldots,\xi_n)=c_i,
  \qquad i=1,\ldots,n,
\]
\uline{où les} \(\sigma_i\) \uline{désignent les fonctions symétriques élémentaires}. \uline{Soient} \(p_1,\ldots,p_k\) \uline{des entiers} \(\geq1\), \uline{de somme} \(n\), \uline{et pour} \(1\le j\le k\),
\[
  \xi^{(j)}=(\xi_1^{(j)},\ldots,\xi_{p_j}^{(j)}),
\]
\uline{avec}
\[
  \xi_i^{(j)}=\xi_{\sum_{\alpha<j}p_\alpha+i}.
\]
\uline{Posons}~:
\[
  c_i^{(j)}=\sigma_i(\xi_1^{(j)},\ldots,\xi_{p_j}^{(j)}),
  \qquad 1\le j\le k,
  \quad 1\le i\le p_j,
\]
\uline{et soit} \(B\) \uline{la sous-algèbre de} \(C\) \uline{engendrée par les} \(c_i^{(j)}\). \uline{Alors} :
\begin{enumerate}[label=\textup{\roman*)}]
\item \uline{les éléments de la forme}~:
\[
  (\xi^{(1)})^{\alpha(1)}\ldots (\xi^{(k)})^{\alpha(k)},
\]
\uline{où} \(\alpha(j)=(\alpha_1^{(j)},\ldots,\alpha_{p_j}^{(j)})\), et \(0\le \alpha_i^{(j)}\le p_j-i\), \uline{forment une base de} \(C\) \uline{considéré comme} \(B\)-\uline{module}.
\item \uline{les générateurs} \(c_i^{(j)}\) \uline{de la} \(A\)-\uline{algèbre} \(B\) \uline{satisfont aux} \(n\) \uline{relations isobares déduites de la relation}~::
\[
  \sum_{i=1}^{n}c_i=\prod_{j=1}^{k}\left(\sum_{i=1}^{p_j}c_i^{(j)}\right)
\]
